\chapter{Evaluation Model} %Abstract
Expressions in EDADSL are evaluated according to standard mathematical precedence rules. The evaluation model is eager (strict), meaning all sub-expressions are evaluated before the operator is applied, with the following exceptions:

\begin{itemize}
\item \texttt{\&\&} and \texttt{||}: Short-circuit evaluation. \texttt{a \&\& b} evaluates \texttt{b} only if \texttt{a} is \texttt{True}. \texttt{a || b} evaluates \texttt{b} only if \texttt{a} is \texttt{False}.
\item \texttt{?->}: Conditional expressions only evaluate the taken branch.
\item \texttt{*}: If either operand evaluates to \texttt{0} (int or real), the result is \texttt{0} and the other operand is not evaluated.
\end{itemize}

\section{Evaluation Order}
Within an expression, operands are evaluated left-to-right. Function arguments are evaluated left-to-right before the function is called.

\section{Operator Precedence}
% Reference to Chapter 7 (Expressions)
The complete operator precedence is defined in Section~\ref{sec:operators}.

\section{Type Coercion}
Type coercion follows the hierarchy: \texttt{int} $<$ \texttt{real} $<$ \texttt{complex}.
\begin{itemize}
\item \texttt{int} coerces to \texttt{real}: \texttt{5} becomes \texttt{5.0}
\item \texttt{real} coerces to \texttt{complex}: \texttt{3.14} becomes \texttt{3.14 + 0j}
\item \texttt{int} coerces to \texttt{complex}: \texttt{5} becomes \texttt{5 + 0j}
\item \texttt{bool} coerces to \texttt{int}: \texttt{True} becomes \texttt{1}
\item \texttt{bool} coerces to \texttt{real}: \texttt{True} becomes \texttt{1.0}
\item \texttt{bool} coerces to \texttt{complex}: \texttt{True} becomes \texttt{1 + 0j}
\item \texttt{bool} coerces to \texttt{string}: \texttt{True} becomes \texttt{"true"}, \texttt{False} becomes \texttt{"false"}
\item Mixed numeric operations promote to the wider type: \texttt{int + real} $\rightarrow$ \texttt{real}, \texttt{real + complex} $\rightarrow$ \texttt{complex}
\item Coercion is never lossy in the upward direction
\item Downward coercion (\texttt{complex} $\rightarrow$ \texttt{real}, \texttt{real} $\rightarrow$ \texttt{int}) raises a \texttt{Type Error} unless done explicitly via a casting function (e.g., \texttt{f.int()}, \texttt{f.real()})
\end{itemize}


\chapter{Execution Order}  %Concrete

\section{Sequential Execution}
Program statements are executed sequentially from the entry point (\texttt{start}). Electrical statements (\texttt{elec.part}, \texttt{elec.connect}, \texttt{elec.net}) execute immediately and register components to the electrical graph. Physical constraints (\texttt{phys.*}) execute immediately but their resolution is deferred until after program execution completes.

\section{Expression Evaluation}
Within a single expression, operands are evaluated left-to-right. The pipeline operator \texttt{|=>} evaluates the left operand first, then passes the result to the function reference on the right. Function arguments are evaluated left-to-right before the function is called. Short-circuit operators (\texttt{\&\&}, \texttt{||}) may skip evaluation of the right operand when the result is determined, unless the unevaluated operand contains a function call that modifies global scope, in which case both operands are always evaluated.

\section{Assertions}
The \texttt{assert} statement evaluates its boolean condition at runtime. If the condition is \texttt{True}, execution continues. If \texttt{False}, the specified error is raised immediately with the provided message. The function form \texttt{f.util.assert()} behaves identically. Assertions are not optimized away at compile time.

\section{Error Handling}
The \texttt{try/except} construct provides structured error handling. When an error occurs inside a \texttt{try} block, the runtime searches the \texttt{otherwise} clauses in order for a matching error type. If found, the error is caught and the corresponding block executes. If not found, the error propagates after the \texttt{except} block (if present). The \texttt{except} block always executes, whether the \texttt{try} succeeded, an error was caught, or an unhandled error propagated.

\section{Function Value Calling}
A function value stored in a variable is called using the syntax \texttt{\$var(args)}. The evaluation proceeds as follows:
\begin{enumerate}
\item The variable is evaluated to obtain the function value
\item If the variable holds \texttt{null}, a \texttt{Runtime Error} is raised
\item If the variable does not hold a \texttt{func} value, a \texttt{Type Error} is raised
\item All arguments are evaluated left-to-right
\item The function is called with the evaluated arguments
\item If the function does not accept the provided argument count or types, a \texttt{Type Error} is raised
\end{enumerate}

\section{Safe Function Calling}
The \texttt{?()} and \texttt{??()} operators provide null-safe function invocation. Both evaluate all arguments \textbf{before} checking the function reference. Errors in argument evaluation are never suppressed.

\subsection{Safe Call (\texttt{?()})}
\begin{enumerate}
\item All arguments are evaluated left-to-right
\item The variable is evaluated to obtain the function value
\item If the variable holds \texttt{null}: returns \texttt{null} (no error)
\item If the variable does not hold a \texttt{func} value: raises a \texttt{Type Error}
\item The function is called with the evaluated arguments
\item Errors raised within the function \textbf{propagate} normally
\end{enumerate}

\subsection{Null-Safe Call (\texttt{??()})}
\begin{enumerate}
\item All arguments are evaluated left-to-right
\item The variable is evaluated to obtain the function value
\item If the variable holds \texttt{null}: returns \texttt{null} (no error)
\item If the variable does not hold a \texttt{func} value: returns \texttt{null} (no error)
\item The function is called with the evaluated arguments
\item If the function call raises \textbf{any} error (\texttt{Type}, \texttt{Value}, \texttt{Division by Zero}, \texttt{Runtime}, \texttt{Assertion}): returns \texttt{null} (error suppressed)
\end{enumerate}
The \texttt{??()} operator is equivalent to wrapping the call in \texttt{try \{ fn(args) \} otherwise \{ null \}}, but without introducing a new code block.

\section{Function Composition}
The composition operator \texttt{<o} creates a new function value at evaluation time. The right operand is evaluated first, then the left operand. Both must evaluate to \texttt{func} values. The result is a new anonymous function that, when called, evaluates \texttt{\$g(x)} first, then passes the result to \texttt{\$f}. The composed function captures both function values by reference at composition time.

\section{Declaration Order}
Variables must be declared before they are accessed. Reading a variable before its declaration statement has been executed raises a \texttt{Runtime Error}. Declarations are processed in order during execution (there is no separate declaration phase). Variable types may change on redeclaration.


\chapter{Constraint Resolution}
Physical constraints are resolved by the solver according to the priority system defined in Section~\ref{sec:physical-constraints}.

\section{Conflict Resolution}
When constraints conflict, the solver applies the higher-priority constraint and optimizes the lower-priority one. See Section~\ref{sec:solver-model} for the complete resolution rules.


\chapter{Type System Rules}
EDADSL uses a dynamic type system with distinct types for electrical and non-electrical objects.

\section{Type Categories}
\begin{itemize}
\item Generic types: Boolean, Integer, Real, Complex, String, List, Set, Function
\item Electrical types: Part, Pin, Connection, Net
\end{itemize}

\section{Type Safety}
EDADSL performs type checking at both compile time and runtime:

\begin{itemize}
\item \textbf{Compile-time:} Type mismatches in variable declarations (e.g., \texttt{\$x [int]= "hello"}) are detected when the script is parsed.
\item \textbf{Runtime:} Type mismatches in function arguments and operator usage are detected when the code is executed. A \texttt{Type Error} is raised for incompatible types.
\item \textbf{Sets:} Sets can contain any data type. Mixed-type sets are allowed.
\item \textbf{Coercion:} Numeric types coerce upward: \texttt{int} $\to$ \texttt{real} $\to$ \texttt{complex}. Other types do not coerce. Function values do not coerce to or from any other type.
\end{itemize}

\section{Error Types}
EDADSL defines seven error types. Each error type has a full name and a short alias for use in \texttt{assert} statements and \texttt{f.util.assert} calls.

\begin{center}
\begin{tabular}{lll}
Full Name & Alias & Typical Cause \\ \hline
\texttt{Type Error} & \texttt{terr} & Incompatible types \\
\texttt{Value Error} & \texttt{verr} & Invalid value (e.g., out of domain) \\
\texttt{Division by Zero Error} & \texttt{0div} & Division by zero \\
\texttt{Runtime Error} & \texttt{rerr} & Invalid runtime state \\
\texttt{Assertion Error} & \texttt{aerr} & Failed assertion \\
\texttt{Permission Error} & \texttt{perm} & Filesystem sandbox violation \\
\end{tabular}
\end{center}

All errors halt program execution immediately. There is no try/catch mechanism.\part{Core Formal Model}
